% Estilo del Manual de Lázaro Linux.
% Aquí solo va la apariencia; el texto está en capitulos/.

% --- Idioma y letras -------------------------------------------------------
\usepackage[T1]{fontenc}
\usepackage[utf8]{inputenc}
\usepackage[spanish,es-noshorthands,es-tabla]{babel}
\usepackage{tgpagella}            % texto: Pagella, fácil de leer en pantalla y papel
\usepackage{tgheros}              % títulos y recuadros: Heros (sin remates)
\renewcommand{\ttdefault}{lmtt}  % comandos de la terminal: Latin Modern Mono
\usepackage{microtype}
\usepackage{amssymb}
\linespread{1.08}

% --- Página ---------------------------------------------------------------
\usepackage[a4paper,top=2.4cm,bottom=2.6cm,inner=2.6cm,outer=2.4cm,headheight=15pt]{geometry}
\usepackage[dvipsnames]{xcolor}
\usepackage{graphicx}
\graphicspath{{imagenes/}{imagenes/capturas/}{../live/usr/share/calamares/branding/lazaro/}}
\usepackage{booktabs,tabularx,array}
\usepackage{enumitem}
\setlist{itemsep=0.25em}
\usepackage{fancyhdr}
\usepackage{titlesec}
\usepackage{tcolorbox}
\tcbuselibrary{skins,breakable}
\usepackage{csquotes}
\usepackage{menukeys}             % \menu{Archivo > Guardar}, \keys{\ctrl + C}
\usepackage[backend=biber,style=numeric,sorting=none,refsegment=chapter,defernumbers=true]{biblatex}
\addbibresource{referencias.bib}
\usepackage{hyperref}

% --- Colores de Lázaro ----------------------------------------------------
\definecolor{pulso}{HTML}{C62828}     % rojo del logo (el pulso)
\definecolor{tinta}{HTML}{263238}
\definecolor{consejo}{HTML}{2E7D32}
\definecolor{cuidado}{HTML}{E65100}
\definecolor{nota}{HTML}{1565C0}

\hypersetup{
  colorlinks=true, linkcolor=tinta, urlcolor=nota, citecolor=consejo,
  pdftitle={Manual de Lázaro Linux},
  pdfauthor={Proyecto Lázaro Linux},
  pdfsubject={Guía de uso de Lázaro Linux para todos},
  pdflang={es},
  bookmarksnumbered=true,
}

% --- Títulos ----------------------------------------------------------------
\titleformat{\chapter}[display]
  {\sffamily\bfseries\color{tinta}}
  {\color{pulso}\Large Capítulo \thechapter}
  {0.4em}
  {\Huge}
  [\vspace{0.3em}{\color{pulso}\titlerule[1.5pt]}]
\titlespacing*{\chapter}{0pt}{-1.5em}{1.8em}
\titleformat{\section}{\sffamily\bfseries\Large\color{tinta}}{\color{pulso}\thesection}{0.8em}{}
\titleformat{\subsection}{\sffamily\bfseries\large\color{tinta}}{\thesubsection}{0.8em}{}
\titleformat{\part}[display]{\centering\sffamily\bfseries\color{tinta}}
  {\color{pulso}\huge Parte \thepart}{0.6em}{\Huge}

\pagestyle{fancy}
\fancyhf{}
\renewcommand{\headrulewidth}{0.4pt}
\fancyhead[LE]{\sffamily\small\leftmark}
\fancyhead[RO]{\sffamily\small\rightmark}
\fancyfoot[C]{\sffamily\small\thepage}
\fancypagestyle{plain}{\fancyhf{}\renewcommand{\headrulewidth}{0pt}\fancyfoot[C]{\sffamily\small\thepage}}

% --- Recuadros --------------------------------------------------------------
% \begin{consejo} ... \end{consejo}   algo que ayuda
% \begin{cuidado} ... \end{cuidado}   algo que puede salir mal
% \begin{nota} ... \end{nota}         un dato extra
\newcommand{\recuadro}[3]{%
  \newtcolorbox{#1}{enhanced,breakable,colback=#2!6,colframe=#2,
    boxrule=0pt,leftrule=3pt,arc=0pt,outer arc=0pt,
    left=8pt,right=8pt,top=6pt,bottom=6pt,
    fonttitle=\sffamily\bfseries\color{#2},coltitle=#2,colbacktitle=#2!6,
    title={#3},attach title to upper={\par\smallskip}}}
\recuadro{consejo}{consejo}{Consejo}
\recuadro{cuidado}{cuidado}{Cuidado}
\recuadro{nota}{nota}{Para saber más}

% Pasos numerados grandes: \begin{pasos} \item ... \end{pasos}
\newcommand{\numpaso}[1]{\tcbox[on line,colback=pulso,colframe=pulso,coltext=white,
  boxsep=1pt,left=3pt,right=3pt,top=1pt,bottom=1pt,arc=2pt]{\sffamily\bfseries #1}}
\newlist{pasos}{enumerate}{1}
\setlist[pasos]{label=\protect\numpaso{\arabic*},leftmargin=2.2em,itemsep=0.5em}

% Lo que se escribe en la terminal.
\newtcolorbox{terminal}{enhanced,colback=tinta,colupper=white,boxrule=0pt,arc=2pt,
  left=8pt,right=8pt,top=5pt,bottom=5pt,fontupper=\ttfamily\footnotesize}
\newcommand{\comando}[1]{\texttt{#1}}

% --- Capturas de pantalla ----------------------------------------------------
% \captura{archivo}{pie de foto}
% Si la imagen todavía no existe, deja un recuadro gris para no frenar la
% escritura: el PDF se compila igual y se ve qué capturas faltan.
\newcommand{\captura}[3][0.85\textwidth]{%
  \begin{figure}[htbp]\centering
  \IfFileExists{imagenes/capturas/#2.png}{%
    \fcolorbox{black!25}{white}{\includegraphics[width=#1]{#2}}}{%
    \IfFileExists{imagenes/capturas/#2.jpg}{%
      \fcolorbox{black!25}{white}{\includegraphics[width=#1]{#2}}}{%
      \fcolorbox{black!30}{black!5}{\parbox[c][0.28\textheight][c]{#1}{\centering
        \sffamily\color{black!55}Captura pendiente:\\\texttt{#2}}}}}
  \caption{#3}\label{fig:#2}
  \end{figure}}

% Términos del glosario: \termino{UEFI} pone la palabra en cursiva la
% primera vez; el glosario completo está en el apéndice.
\newcommand{\termino}[1]{\emph{#1}}

% Referencias al final de cada capítulo: \fuentes imprime solo las citadas
% en ese capítulo. La lista completa va en el apéndice.
\newcommand{\fuentes}{\printbibliography[segment=\therefsegment,
  heading=subbibliography,title={Fuentes y lecturas del capítulo}]}
